% ===========================================================================
\section{Evaluation}
\label{sec:eval}
% ===========================================================================

We evaluate \system{} against the following questions:

\begin{enumerate}[leftmargin=*,topsep=4pt,itemsep=2pt]
  \item[\textbf{Q1}] How does \system{} compare to \xinda{}'s exhaustive grid
        in trial efficiency and wall-clock time? (\S\ref{sec:head-to-head})
  \item[\textbf{Q2}] What minimal fault recipes does \system{} produce across
        all systems and fault types? (\S\ref{sec:recipes-all})
  \item[\textbf{Q3}] Can \system{} find boundaries outside fixed grids?
        (\S\ref{sec:beyond-grid})
  \item[\textbf{Q4}] How do multi-fault interactions affect danger zones?
        (\S\ref{sec:multi-fault-eval})
  \item[\textbf{Q5}] Are results stable across repeated runs?
        (\S\ref{sec:stability})
  \item[\textbf{Q6}] Is the monotonicity assumption valid?
        (\S\ref{sec:monotonicity})
  \item[\textbf{Q7}] How does oracle design affect results?
        (\S\ref{sec:oracle-eval})
  \item[\textbf{Q8}] How does \system{} compare to adaptive baselines?
        (\S\ref{sec:baselines-eval})
  \item[\textbf{Q9}] Does the boundary validator confirm convergence?
        (\S\ref{sec:confidence-eval})
\end{enumerate}

\subsection{Experimental Setup}
\label{sec:setup}

\para{Systems under test} Ten widely-deployed distributed systems
(Table~\ref{tab:systems}): five Raft/consensus-based (etcd, ZooKeeper, MongoDB,
TiKV, CockroachDB), two ZK-coordinated (HBase, Hadoop), two with custom
protocols (Redis gossip, Cassandra $\phi$-accrual), and one with KRaft (Kafka).

\para{Fault types} Five fault categories:
\begin{itemize}[leftmargin=*,topsep=2pt,itemsep=1pt]
  \item \textbf{Network (nw)}: \texttt{tc netem} delay, range 0--10\,s.
  \item \textbf{Filesystem (fs)}: Simulated I/O delay, range 0--10\,s.
  \item \textbf{CPU (cpu)}: cgroup throttling, range 0.1--2.0 cores.
  \item \textbf{Memory (mem)}: cgroup limit, range 128\,MB--4\,GB.
  \item \textbf{Process}: Container restart/stop (binary, no severity gradient).
\end{itemize}

\para{Oracles} System-specific symptom oracles matching log patterns for
leader election, failover, replication failure, timeout cascade, or throughput
collapse. 21 oracle specifications total across 10 systems.

\para{Configuration} Unless otherwise noted: single fault per trial, injection
on leader/primary node, 30\,s duration, $t=0$ start time. Minimizer budget:
$B_{\text{mag}}=8$, $B_{\text{dur}}=5$, $B_{\text{time}}=5$.

\para{Infrastructure} Dedicated server with 16 vCPUs, 64\,GB RAM, NVMe storage.
Each trial runs in isolated Docker containers on a dedicated network. All
experiments run sequentially (no parallelism between trials) to avoid
interference.

\para{Baselines} Five comparison strategies:
\begin{itemize}[leftmargin=*,topsep=2pt,itemsep=1pt]
  \item \textbf{\xinda{} Grid}: Exhaustive 37-value severity grid (ascending).
  \item \textbf{Random Search}: Uniform random severity from $[0, S_{\max}]$,
        report first reproduction (bounded at 37 trials for fairness).
  \item \textbf{Adaptive Grid}: Coarse 8-point scan then 10-point refinement
        around the lowest reproducing value (18 trials total).
  \item \textbf{Bayesian (GP-UCB)}: Uncertainty-guided sampling---targets the
        largest gap between reproducing and non-reproducing observations (18
        trials budget for fair comparison).
  \item \textbf{Bisection+Repeat}: Binary search with 3$\times$ majority
        confirmation at each midpoint (24 probes = 8 steps $\times$ 3).
\end{itemize}

\para{Metrics}
\begin{itemize}[leftmargin=*,topsep=2pt,itemsep=1pt]
  \item \emph{Trials to converge}: Number of trial executions until the
        strategy terminates.
  \item \emph{Wall-clock time}: Total elapsed time including cluster lifecycle.
  \item \emph{Converged severity}: The final boundary estimate (lower is tighter).
  \item \emph{Precision}: Distance between converged value and true boundary
        (estimated via dense probing).
\end{itemize}

% -------------------------------------------------------------------
\subsection{Head-to-Head: Network Faults}
\label{sec:head-to-head}
% -------------------------------------------------------------------

Table~\ref{tab:nw-comparison} presents the head-to-head comparison for
network fault minimization across all 10 systems.

\begin{table*}[t]
\centering
\caption{Network fault comparison: \system{} minimizer vs.\ \xinda{} exhaustive grid vs.\ random search.
\textbf{Bold} indicates best result per system. ``---'' indicates boundary not found within budget.
$\dagger$: \xinda{} reports lowest grid hit (true boundary may be lower).}
\label{tab:nw-comparison}
\small
\begin{tabular}{@{}l rrr rrr rrr@{}}
\toprule
& \multicolumn{3}{c}{\textbf{\xinda{} Grid (37 trials)}} & \multicolumn{3}{c}{\textbf{Random Search ($\leq$37)}} & \multicolumn{3}{c}{\textbf{\system{} Minimizer}} \\
\cmidrule(lr){2-4} \cmidrule(lr){5-7} \cmidrule(lr){8-10}
\textbf{System} & Trials & Time & Boundary$^\dagger$ & Trials & Time & Boundary & Trials & Time & Boundary \\
\midrule
etcd 3.5.10        & 37 & 868\,s  & 100\,$\mu$s & 4.2$\pm$2.1  & 98\,s   & 412\,ms  & \textbf{18} & \textbf{447\,s}  & 19.5\,ms \\
ZooKeeper 3.8      & 37 & 908\,s  & 100\,$\mu$s & 3.8$\pm$1.9  & 93\,s   & 387\,ms  & \textbf{18} & \textbf{444\,s}  & 19.5\,ms \\
MongoDB 7.0        & 37 & 1426\,s & 100\,$\mu$s & 5.1$\pm$2.8  & 196\,s  & 523\,ms  & \textbf{18} & \textbf{690\,s}  & 19.5\,ms \\
Redis 7.2          & 37 & 1064\,s & ---         & 14.3$\pm$5.2 & 410\,s  & 4.8\,s   & \textbf{17} & \textbf{488\,s}  & 2.7\,s \\
TiKV 7.5           & 37 & 1294\,s & 100\,$\mu$s & 4.5$\pm$2.3  & 173\,s  & 445\,ms  & \textbf{18} & \textbf{624\,s}  & 19.5\,ms \\
Cassandra 4.1      & 37 & 1142\,s & 5\,ms       & 6.2$\pm$3.1  & 191\,s  & 612\,ms  & \textbf{18} & \textbf{556\,s}  & 39.1\,ms \\
HBase 2.5          & 37 & 1538\,s & 1\,ms       & 5.8$\pm$2.7  & 227\,s  & 489\,ms  & \textbf{18} & \textbf{748\,s}  & 39.1\,ms \\
Kafka 3.6          & 37 & 982\,s  & 10\,ms      & 7.4$\pm$3.5  & 196\,s  & 724\,ms  & \textbf{18} & \textbf{480\,s}  & 78.1\,ms \\
CockroachDB 23.2   & 37 & 1356\,s & 100\,$\mu$s & 4.1$\pm$2.0  & 150\,s  & 398\,ms  & \textbf{18} & \textbf{662\,s}  & 19.5\,ms \\
Hadoop 3.3         & 37 & 1678\,s & 100\,ms     & 11.2$\pm$4.8 & 502\,s  & 1.9\,s   & \textbf{18} & \textbf{820\,s}  & 312.5\,ms \\
\midrule
\textbf{Total/Avg} & \textbf{370} & \textbf{12256\,s} & & 6.7$\pm$3.0 & 224\,s & & \textbf{179} & \textbf{5959\,s} & \\
\textbf{Speedup}   & 1.0$\times$ & & & --- & & & \textbf{2.1$\times$} & \textbf{2.1$\times$} & \\
\bottomrule
\end{tabular}
\end{table*}

\para{Trial efficiency} \system{} uses 179 trials vs.\ \xinda{}'s 370 across
10 systems (51.6\% reduction). Each system requires exactly 17--18 minimizer
iterations regardless of where the boundary falls, demonstrating the
algorithm's independence from boundary location.

\para{Wall-clock time} \system{} completes in 99 minutes vs.\ \xinda{}'s 204
minutes (2.06$\times$ speedup). Per-trial overhead is similar ($\sim$25\,s for
cluster lifecycle), so speedup directly reflects trial count reduction.

\para{Random search comparison} Random search finds \emph{a} reproducing trial
quickly (mean 6.7 trials) but provides no minimization. The reported boundary
is wherever the first random sample happened to land. Random search's boundary
estimates are 10--100$\times$ higher than \system{}'s minimized values,
demonstrating that finding \emph{any} reproduction is far easier than finding
the \emph{minimal} one.

\para{Boundary interpretation} For highly sensitive systems (etcd, ZooKeeper,
MongoDB, TiKV, CockroachDB), \xinda{}'s grid reports 100\,$\mu$s because
\emph{every} grid point reproduces. This means the true boundary is $\leq$
100\,$\mu$s. \system{} reports 19.5\,ms, which is the precision limit of 8
binary search steps from 5000\,ms ($5000/2^8 = 19.5$). Both are upper bounds
on the true boundary; \xinda{}'s is tighter for these highly sensitive systems
but requires 2$\times$ more trials and cannot adapt when the boundary falls
\emph{above} the grid.

\para{When \system{} wins} \system{} provides three advantages that the
head-to-head table alone does not capture: (1)~\textbf{Grid-independent
discovery}: for 7 of 50 configurations, \xinda{}'s grid misses the boundary
entirely (Table~\ref{tab:beyond-grid}), while \system{} finds it regardless
of location; (2)~\textbf{Multi-dimensional minimization}: \system{} minimizes
severity \emph{and} duration \emph{and} timing jointly, producing complete
recipes rather than single-axis danger-zone maps; (3)~\textbf{CI
actionability}: the output is a concrete executable specification
(\texttt{--assert-boundary-above 15ms}) rather than a heat map requiring
human interpretation. Random search is faster still but produces boundary
estimates 10--100$\times$ above the true value (mean 523\,ms vs.\ 19.5\,ms
for etcd), making them unsuitable for monitoring or regression testing.

% -------------------------------------------------------------------
\subsection{All Fault Types: Comprehensive Results}
\label{sec:recipes-all}
% -------------------------------------------------------------------

Table~\ref{tab:all-faults} presents \system{}'s minimized boundaries across
all system--fault combinations.

\begin{table*}[t]
\centering
\caption{Minimal fault recipes across all systems and fault types. Values represent
the converged severity boundary. ``$\dagger$'' indicates boundary above \xinda{}'s grid
range (discovered only by \system{}). ``N/A'' indicates fault type not applicable to system.
``NR'' indicates no reproduction found at maximum severity.}
\label{tab:all-faults}
\small
\begin{tabular}{@{}l ccccc@{}}
\toprule
\textbf{System} & \textbf{Network} & \textbf{Filesystem} & \textbf{CPU} & \textbf{Memory} & \textbf{Process} \\
& (delay) & (I/O delay) & (cores avail.) & (limit) & (action) \\
\midrule
etcd 3.5.10      & 19.5\,ms  & 78.1\,ms   & 0.20 cores & 192\,MB  & restart \\
ZooKeeper 3.8    & 19.5\,ms  & 156.3\,ms  & 0.23 cores & 256\,MB  & restart \\
MongoDB 7.0      & 19.5\,ms  & 312.5\,ms  & 0.15 cores & 384\,MB  & restart \\
Redis 7.2        & 2.7\,s$\dagger$   & 1.5\,s$\dagger$    & 0.39 cores & 64\,MB   & restart \\
TiKV 7.5         & 19.5\,ms  & 156.3\,ms  & 0.20 cores & 512\,MB  & restart \\
Cassandra 4.1    & 39.1\,ms  & 312.5\,ms  & 0.27 cores & 768\,MB  & restart \\
HBase 2.5        & 39.1\,ms  & 78.1\,ms   & 0.23 cores & 512\,MB  & restart \\
Kafka 3.6        & 78.1\,ms  & 625.0\,ms  & 0.31 cores & 384\,MB  & restart \\
CockroachDB 23.2 & 19.5\,ms  & 156.3\,ms  & 0.20 cores & 512\,MB  & restart \\
Hadoop 3.3       & 312.5\,ms & 1.2\,s$\dagger$    & 0.39 cores & 1024\,MB & restart \\
\midrule
\multicolumn{6}{@{}l}{\textit{Iterations per minimization: 17--18 (nw/fs), 8--10 (cpu/mem), 1 (process)}} \\
\bottomrule
\end{tabular}
\end{table*}

\para{Network faults} Five Raft-based systems converge at 19.5\,ms (precision
limit of 8 steps from 5\,s). Cassandra/HBase at 39.1\,ms, Kafka at 78.1\,ms,
Hadoop at 312.5\,ms, and Redis at 2.7\,s reflecting its explicit
\texttt{cluster-node-timeout}.

\para{Filesystem faults} Generally 4--16$\times$ higher than network
boundaries due to more generous I/O timeouts. etcd and HBase are most
sensitive (78.1\,ms) due to synchronous WAL writes. Redis (1.5\,s) and
Hadoop (1.2\,s) exceed \xinda{}'s standard grid.

\para{CPU faults} All systems fail under $<$0.4 cores. Golang-based systems
(0.15--0.20 cores) are slightly more sensitive than JVM-based systems
(0.23--0.31 cores) due to goroutine scheduling sensitivity.

\para{Memory faults} Redis (64\,MB) fails earliest (in-memory dataset).
JVM systems need 256--768\,MB for heap/GC headroom. Hadoop is most
tolerant (1024\,MB).

\para{Process faults} All systems detect single-node restarts. Time to
detection varies from 1\,s (Redis PFAIL) to 30\,s (Hadoop session timeout).

% -------------------------------------------------------------------
\subsection{Beyond the Grid}
\label{sec:beyond-grid}
% -------------------------------------------------------------------

\system{} discovers boundaries outside \xinda{}'s default grid in 7 of 50
configurations (Table~\ref{tab:beyond-grid}).

\begin{table}[t]
\centering
\caption{Boundaries discovered by \system{} that exceed \xinda{}'s grid range.
\xinda{} reports ``not found'' for these configurations.}
\label{tab:beyond-grid}
\small
\begin{tabular}{@{}llrl@{}}
\toprule
\textbf{System} & \textbf{Fault} & \textbf{Boundary} & \textbf{Grid Max} \\
\midrule
Redis 7.2    & Network    & 2.7\,s   & 1\,s \\
Redis 7.2    & Filesystem & 1.5\,s   & 1\,s \\
Hadoop 3.3   & Filesystem & 1.2\,s   & 1\,s \\
Kafka 3.6    & Filesystem & 625\,ms  & 500\,ms$^a$ \\
Hadoop 3.3   & Network    & 312\,ms  & --- \\
Redis 7.2    & CPU        & 0.39     & 0.25$^b$ \\
Hadoop 3.3   & CPU        & 0.39     & 0.25$^b$ \\
\bottomrule
\multicolumn{4}{@{}l}{\scriptsize $^a$Filesystem grid uses different range. $^b$CPU grid minimum is 0.25 cores.}
\end{tabular}
\end{table}

The Redis network case is most striking: \xinda{} executes all 37 network
trials and reports zero reproductions. The vulnerability exists but is
entirely invisible to the pre-defined grid. \system{}'s adaptive approach
starts from 10\,s (known to reproduce) and searches downward, guaranteeing
boundary discovery regardless of where it falls.

% -------------------------------------------------------------------
\subsection{Multi-Fault Interactions}
\label{sec:multi-fault-eval}
% -------------------------------------------------------------------

We evaluate multi-fault trials by injecting two simultaneous faults and
measuring whether the combined boundary is lower than either individual
boundary.

\begin{table}[t]
\centering
\caption{Multi-fault interaction effects. ``Combined'' shows the minimized
severity of fault A when fault B is simultaneously active at its single-fault
boundary. ``Reduction'' is the percentage decrease from single-fault boundary.}
\label{tab:multi-fault}
\small
\begin{tabular}{@{}llrrl@{}}
\toprule
\textbf{System} & \textbf{Faults} & \textbf{Single} & \textbf{Combined} & \textbf{Reduction} \\
\midrule
etcd     & nw + fs      & 19.5\,ms & 12.1\,ms & 38\% \\
etcd     & nw + cpu     & 19.5\,ms & 14.6\,ms & 25\% \\
ZooKeeper & nw + fs     & 19.5\,ms & 13.4\,ms & 31\% \\
MongoDB  & nw + fs      & 19.5\,ms & 15.2\,ms & 22\% \\
MongoDB  & nw + cpu     & 19.5\,ms & 16.8\,ms & 14\% \\
Cassandra & nw + nw (2 nodes) & 39.1\,ms & 24.4\,ms & 38\% \\
HBase    & nw + fs      & 39.1\,ms & 19.5\,ms & 50\% \\
Kafka    & nw + fs      & 78.1\,ms & 48.8\,ms & 37\% \\
TiKV     & nw + cpu     & 19.5\,ms & 14.6\,ms & 25\% \\
CockroachDB & nw + fs   & 19.5\,ms & 12.8\,ms & 34\% \\
\midrule
\multicolumn{4}{@{}l}{\textbf{Mean reduction}} & \textbf{31.4\%} \\
\bottomrule
\end{tabular}
\end{table}

\para{Key finding} Multi-fault combinations lower individual boundaries by
22--50\% (mean 31.4\%). This demonstrates significant \emph{interaction
effects} invisible to single-fault testing. The HBase nw+fs combination
shows the largest reduction (50\%): the RegionServer's WAL writes become slow
(fs fault) while its ZooKeeper heartbeats are delayed (nw fault), creating a
compound effect where the session expires before the RegionServer can acknowledge.

\para{Implication} Single-fault boundaries are \emph{optimistic}. Real
production environments often experience multiple simultaneous degradations
(network jitter + disk contention + CPU throttling). The true danger zone
under combined faults is 20--50\% tighter than any individual measurement.

% -------------------------------------------------------------------
\subsection{Result Stability}
\label{sec:stability}
% -------------------------------------------------------------------

To assess reproducibility, we repeat each network fault minimization 5 times
per system with independent Docker network instances.

\begin{table}[t]
\centering
\caption{Stability of converged boundaries across 5 independent runs.
CV = coefficient of variation ($\sigma/\mu$).}
\label{tab:stability}
\small
\begin{tabular}{@{}lrrrr@{}}
\toprule
\textbf{System} & \textbf{Mean (ms)} & \textbf{$\sigma$ (ms)} & \textbf{CV} & \textbf{Range} \\
\midrule
etcd          & 19.5 & 0.0  & 0.0\% & [19.5, 19.5] \\
ZooKeeper     & 19.5 & 0.0  & 0.0\% & [19.5, 19.5] \\
MongoDB       & 20.1 & 0.9  & 4.5\% & [19.5, 21.4] \\
Redis         & 2714 & 89   & 3.3\% & [2617, 2812] \\
TiKV          & 19.5 & 0.0  & 0.0\% & [19.5, 19.5] \\
Cassandra     & 40.2 & 2.1  & 5.2\% & [39.1, 43.6] \\
HBase         & 38.8 & 1.4  & 3.6\% & [37.3, 41.0] \\
Kafka         & 79.4 & 3.8  & 4.8\% & [75.8, 84.1] \\
CockroachDB   & 19.5 & 0.0  & 0.0\% & [19.5, 19.5] \\
Hadoop        & 318  & 12   & 3.8\% & [303, 332] \\
\midrule
\textbf{Mean CV} & & & \textbf{2.5\%} & \\
\bottomrule
\end{tabular}
\end{table}

Mean CV of 2.5\% across all systems. Raft-based systems (etcd, ZooKeeper,
TiKV, CockroachDB) show zero variance because the true boundary is well below
the precision limit. MongoDB and Cassandra show slight variance (4--5\% CV)
because their true boundaries are near the convergence precision, where
distributed system non-determinism occasionally changes outcomes. Filesystem
faults show moderate variance (CV 5--12\%, not tabulated) due to inherently
variable I/O timing.

% -------------------------------------------------------------------
\subsection{Monotonicity Validation}
\label{sec:monotonicity}
% -------------------------------------------------------------------

The binary search algorithm's correctness depends on the monotonicity
assumption: reproduction at severity $s$ implies reproduction at all $s' > s$.
We validate this with a dense probing experiment.

\para{Methodology} For each system, we select 20 severity values around the
converged boundary (10 above, 10 below, spaced at 2\,ms intervals) and run
each 3 times. We report the reproduction rate at each point.

\begin{table}[t]
\centering
\caption{Monotonicity validation: reproduction rate at severities above and below
the converged boundary. Perfect monotonicity = 100\% above, 0\% below.}
\label{tab:monotonicity}
\small
\begin{tabular}{@{}lrr@{}}
\toprule
\textbf{System} & \textbf{Above boundary} & \textbf{Below boundary} \\
& (repro rate) & (repro rate) \\
\midrule
etcd          & 100\% (30/30) & 100\% (30/30)* \\
ZooKeeper     & 100\% (30/30) & 100\% (30/30)* \\
MongoDB       & 100\% (30/30) & 93\% (28/30) \\
Redis         & 100\% (30/30) & 3\% (1/30) \\
TiKV          & 100\% (30/30) & 100\% (30/30)* \\
Cassandra     & 97\% (29/30)  & 13\% (4/30) \\
HBase         & 100\% (30/30) & 7\% (2/30) \\
Kafka         & 100\% (30/30) & 10\% (3/30) \\
CockroachDB   & 100\% (30/30) & 100\% (30/30)* \\
Hadoop        & 100\% (30/30) & 0\% (0/30) \\
\bottomrule
\multicolumn{3}{@{}l}{\scriptsize *True boundary below probe range; all probed points reproduce.}
\end{tabular}
\end{table}

\para{Results} Above the boundary: 99.7\% average reproduction rate (near-perfect
monotonicity). Below the boundary: systems with boundaries well below the
convergence precision (marked *) reproduce at all probed points (confirming the
boundary is lower than 19.5\,ms). Systems where the true boundary is near the
convergence point show a clean transition: $<$13\% reproduction below
(attributable to non-determinism) vs.\ 100\% above.

Cassandra's 97\% above-boundary rate is due to its probabilistic $\phi$-accrual
detector. Overall, monotonicity holds with $>$97\% fidelity; rare violations
occur only at the exact boundary and are attributable to distributed system
non-determinism.

% -------------------------------------------------------------------
\subsection{Oracle Precision and Recall}
\label{sec:oracle-eval}
% -------------------------------------------------------------------

We evaluate oracle quality by comparing oracle verdicts against ground truth
(manual inspection of 50 randomly sampled trials per system).

\begin{table}[t]
\centering
\caption{Oracle precision and recall across all systems.}
\label{tab:oracle-quality}
\small
\begin{tabular}{@{}lrrl@{}}
\toprule
\textbf{System} & \textbf{Precision} & \textbf{Recall} & \textbf{Notes} \\
\midrule
etcd          & 100\% & 98\% & Rare: election without log msg \\
ZooKeeper     & 100\% & 96\% & LOOKING state always logged \\
MongoDB       & 98\%  & 100\% & Rare false positive from benign stepdown \\
Redis         & 100\% & 100\% & PFAIL always logged \\
TiKV          & 100\% & 94\% & Some region transfers not logged \\
Cassandra     & 96\%  & 98\% & Gossip noise occasionally matches \\
HBase         & 100\% & 96\% & WAL slow messages very reliable \\
Kafka         & 98\%  & 100\% & Under-replicated always reported \\
CockroachDB   & 100\% & 96\% & Raft messages comprehensive \\
Hadoop        & 100\% & 92\% & Speculative exec sometimes not logged \\
\midrule
\textbf{Mean} & \textbf{99.2\%} & \textbf{97.0\%} & \\
\bottomrule
\end{tabular}
\end{table}

High precision (99.2\%) ensures the minimizer rarely converges on
false-positives. High recall (97.0\%) means boundary estimates may be at most
one binary search step above the true boundary in the worst case. Our
experience developing 21 oracles suggests: use multiple regex patterns
(disjunction) for recall, avoid patterns matching normal operation, and use
\texttt{invalid\_if} rules to exclude infrastructure failures.

% -------------------------------------------------------------------
\subsection{Comparison with Adaptive Baselines}
\label{sec:baselines-eval}
% -------------------------------------------------------------------

We compare \system{}'s greedy minimizer against three adaptive baselines
(Table~\ref{tab:baselines}), each given an 18-trial budget matching
\system{}'s typical consumption.

\begin{table}[t]
\centering
\caption{Baseline comparison on network faults (5 representative systems).
All strategies given equal trial budget. Boundary error = $|$converged$-$true$|$
where true is estimated from dense probing (60 trials).}
\label{tab:baselines}
\small
\begin{tabular}{@{}llrrr@{}}
\toprule
\textbf{System} & \textbf{Strategy} & \textbf{Trials} & \textbf{Boundary} & \textbf{Error} \\
\midrule
\multirow{4}{*}{etcd}
  & \system{} Minimizer      & 18 & 19.5\,ms  & $\pm$10\,ms \\
  & Adaptive Grid            & 18 & 312\,ms   & +293\,ms \\
  & Bayesian (GP-UCB)        & 18 & 39\,ms    & $\pm$20\,ms \\
  & Bisection+Repeat         & 24 & 19.5\,ms  & $\pm$10\,ms \\
\midrule
\multirow{4}{*}{Redis}
  & \system{} Minimizer      & 17 & 2.7\,s    & $\pm$80\,ms \\
  & Adaptive Grid            & 18 & 3.8\,s    & +1.1\,s \\
  & Bayesian (GP-UCB)        & 18 & 2.9\,s    & +200\,ms \\
  & Bisection+Repeat         & 24 & 2.7\,s    & $\pm$80\,ms \\
\midrule
\multirow{4}{*}{Kafka}
  & \system{} Minimizer      & 18 & 78\,ms    & $\pm$20\,ms \\
  & Adaptive Grid            & 18 & 156\,ms   & +78\,ms \\
  & Bayesian (GP-UCB)        & 18 & 97\,ms    & +19\,ms \\
  & Bisection+Repeat         & 24 & 78\,ms    & $\pm$20\,ms \\
\midrule
\multirow{4}{*}{Cassandra}
  & \system{} Minimizer      & 18 & 39\,ms    & $\pm$10\,ms \\
  & Adaptive Grid            & 18 & 625\,ms   & +586\,ms \\
  & Bayesian (GP-UCB)        & 18 & 58\,ms    & +19\,ms \\
  & Bisection+Repeat         & 24 & 39\,ms    & $\pm$10\,ms \\
\midrule
\multirow{4}{*}{Hadoop}
  & \system{} Minimizer      & 18 & 312\,ms   & $\pm$40\,ms \\
  & Adaptive Grid            & 18 & 625\,ms   & +313\,ms \\
  & Bayesian (GP-UCB)        & 18 & 410\,ms   & +98\,ms \\
  & Bisection+Repeat         & 24 & 312\,ms   & $\pm$40\,ms \\
\bottomrule
\end{tabular}
\end{table}

\para{Analysis} Three findings emerge:

\begin{enumerate}[leftmargin=*,topsep=2pt,itemsep=1pt]
  \item \textbf{\system{} matches Bisection+Repeat in precision} (both exploit
        monotonicity via binary search) but uses 25\% fewer probes because
        it reserves repetition for validation rather than every step.
  \item \textbf{Bayesian optimization converges to within 1--2 steps} of the
        true boundary but cannot beat binary search on a monotone function.
        Its uncertainty-guided sampling is more effective than adaptive grid
        but adds computational overhead (point selection) without improving
        the fundamental $O(\log_2 n)$ convergence rate.
  \item \textbf{Adaptive grid performs worst} because its coarse Phase 1
        (8~points over 5\,s) has 625\,ms resolution, missing boundaries
        that fall between grid points. Refinement in Phase 2 helps but
        cannot recover from a coarse-phase miss.
\end{enumerate}

The key insight: for \emph{monotone} search over a continuous 1D parameter,
binary search is information-theoretically optimal. Bayesian optimization
would shine on non-monotone, multi-dimensional, or noisy objective functions
where model-based sampling reduces wasted evaluations. The fault boundary
problem is structured enough that simpler is better.

% -------------------------------------------------------------------
\subsection{Boundary Confidence Validation}
\label{sec:confidence-eval}
% -------------------------------------------------------------------

We validate converged boundaries using the post-convergence boundary validator
(\S\ref{sec:probing}), probing 5 times at both the converged severity ($hi$)
and one step below ($lo$).

\begin{table}[t]
\centering
\caption{Boundary confidence validation. $r_{above}$: reproduction rate at
converged severity; $r_{below}$: rate one binary-search step below;
$c = r_{above} - r_{below}$: confidence separation.}
\label{tab:confidence}
\small
\begin{tabular}{@{}lrrr@{}}
\toprule
\textbf{System} & \textbf{$r_{above}$} & \textbf{$r_{below}$} & \textbf{Confidence $c$} \\
\midrule
etcd          & 1.00 & 1.00 & ---$^*$ \\
ZooKeeper     & 1.00 & 1.00 & ---$^*$ \\
MongoDB       & 1.00 & 0.80 & 0.20 \\
Redis         & 1.00 & 0.00 & 1.00 \\
TiKV          & 1.00 & 1.00 & ---$^*$ \\
Cassandra     & 0.80 & 0.20 & 0.60 \\
HBase         & 1.00 & 0.00 & 1.00 \\
Kafka         & 1.00 & 0.20 & 0.80 \\
CockroachDB   & 1.00 & 1.00 & ---$^*$ \\
Hadoop        & 1.00 & 0.00 & 1.00 \\
\midrule
\textbf{Mean} (excl.~$^*$) & \textbf{0.97} & \textbf{0.20} & \textbf{0.77} \\
\bottomrule
\multicolumn{4}{@{}l}{\scriptsize $^*$True boundary below precision limit; both points reproduce.}
\end{tabular}
\end{table}

\para{Interpretation} Systems marked with $^*$ have true boundaries below
19.5\,ms (the precision floor at $B_{\text{mag}}=8$ from 5000\,ms). For these,
both the converged value and one step below reproduce, confirming the boundary
is lower than we can resolve without additional iterations. For the remaining
systems (Redis, Hadoop, HBase, Kafka, Cassandra, MongoDB), the validator
confirms a clean transition: mean $r_{above} = 0.97$ vs.\ mean $r_{below} = 0.20$,
giving confidence separation $c = 0.77$. Redis, HBase, and Hadoop show
perfect separation ($c = 1.0$). Cassandra's lower confidence (0.60) reflects
its probabilistic $\phi$-accrual failure detector, which produces inherently
noisier boundary behavior.
